\scanpage{186}
% L01–02
ou encore
\[
\rho(a)\deleted{\unsure{\rho}}(a)\in\rho(\widetilde U)\deleted{\unsure{g^{-1}|\rho^{-1}}}.\deleted{\unsure{g}}\widetilde U\bigr)
\]

% L03–07
Posons $\widetilde U_0=\widetilde U=G_{0\,s_0}$, $\underset{\substack{\shortparallel\\\rho(\widetilde U)}}{\widetilde U_1}=G_{0\,t_0}$, la condition envisagée s'écrit donc
\[
\rho(a)\,a\in\widetilde U_1\widetilde U_0\tag*{(24)}
\]
(en plus de la relation
\[
a\in\operatorname{Norm}_{G_0}(\widetilde U_1)\bigr).
\]

% L08–14
\uline{NB} Dans le cas où le polyèdre $P_0$ est orienté, il est plus sympathique dans cette \unsure{traduction} de travailler dans $G_0^+$ (qui est encore transitif sur $S$ et $=$ sur l'ens. des arêtes orientées, qui est $=$ un \uline{torseur} sous $\unsure{G_\infty^+|G_0^+}$), en décrivant $a$ par un élément de $G_0^+$
\[
a\in\operatorname{Norm}_{G_0^+}(\overset{\deleted{\sim}}{U}_0)\qquad U_0=G^+_{0\,s_0}=\widetilde U_0\cap G_0^+ \tag*{(25)}
\]
et en écrivant \deleted{\unsure{(2|la}} alors (24) \inserted{dans $G_0^+$,} sous la forme équivalente
\[
\rho(a)a\in\overset{\deleted{\sim}}{U}_1.U_0\ .\tag*{(24')}
\]

% L15–20
On dit que l'homothétie \inserted{$\lambda$} est stricte si elle satisfait la condition que la structure de graphe \inserted{via} $A_\lambda\in\unsure{\mathfrak P_2(S)|\mathrm P_2(S)}$ définie par
\[
\deleted{\{s,t\}\in}\left\{\begin{aligned}
&A_\lambda=\bigl\{\{s,\lambda t\}\in\mathfrak P_2(S)\mid(s,t)\in S^2,\ \{s,t\}\in A_0\bigr\}\\
&\text{i.e. }\deleted{\{s,t\}\in\unsure{A}}\ \forall s,t\in A_0\qquad
\{s,t\}\in A_0\Longleftrightarrow\{s,\lambda t\}\in A_\lambda
\end{aligned}\right.\tag*{(25)}
\]
est isomorphe à $(S,A_{\unsure{0|\lambda}})$, \deleted{\illegible{1}} condition qu'on

% archive sheet number: 185
